\begin{proofbox}
	\textbf{\textcolor{CProof}{思路提示}}

	要得到 $P(A_i\mid B)$，先按条件概率写为 $\frac{P(A_i\cap B)}{P(B)}$，再分别处理分子和分母：分子用乘法公式，分母用全概率公式。

	\medskip
	\textbf{\textcolor{CProof}{正式推导}}
	\begin{description}[style=nextline, leftmargin=2.8em, labelsep=0.5em, itemsep=0.55em]
		\item[\textbf{步骤一（条件概率）：}] 由条件概率定义（$P(B)>0$ 保证分母非零），
			\[
				P(A_i\mid B) = \frac{P(A_i\cap B)}{P(B)}.
			\]
			\par\textbf{理由：} 条件概率的定义是基础公式 $P(A\mid B)=\frac{P(AB)}{P(B)}$。

		\item[\textbf{步骤二（分子转化）：}] 由乘法公式，$A_i$ 与 $B$ 同时发生的概率可表示为先验概率与似然概率之积，
			\[
				P(A_i\cap B) = P(A_i)P(B\mid A_i).
			\]
			\par\textbf{理由：} 乘法公式 $P(AB)=P(A)P(B\mid A)$ 适用于任何事件，这里 $P(A_i)>0$ 保证条件概率有意义。

		\item[\textbf{步骤三（分母全概率）：}] 由于 $A_1,\dots,A_n$ 是样本空间的一个划分，事件 $B$ 可分解为互斥的 $A_i\cap B$ 的和，故由全概率公式，
			\[
				\begin{aligned}
					P(B) &= \sum_{j=1}^{n} P(A_j\cap B) \\
					&= \sum_{j=1}^{n} P(A_j)P(B\mid A_j).
			\end{aligned}
			\]
			\par\textbf{理由：} 全概率公式的前提是事件组构成划分，且各 $P(A_j)>0$；这里条件已满足。

		\item[\textbf{步骤四（代入合并）：}] 将步骤二和步骤三的结果代入步骤一，
			\[
				P(A_i\mid B) = \frac{P(A_i)P(B\mid A_i)}{\displaystyle\sum_{j=1}^{n} P(A_j)P(B\mid A_j)}.
			\]
			\par\textbf{理由：} 分子分母同时代入，即得到贝叶斯公式的标准形式。
	\end{description}

	\medskip
	\textbf{\textcolor{CProof}{结论回扣}}

	\textbf{因此，后验概率 $P(A_i\mid B)$ 可以通过先验概率 $P(A_i)$ 和似然概率 $P(B\mid A_i)$ 按上述公式计算。需要注意的是，分母是各原因产生 $B$ 的概率之和，保证了 $\sum_{i=1}^{n}P(A_i\mid B)=1$。}
\end{proofbox}
